\chapter{AI Components}
\label{ch:ai}

This chapter details the two AI tracks of \padonap{}: Track~A, an
XGBoost seven-class detector; and Track~B, a Transformer+LSTM
four-horizon forecaster. Both tracks were trained on CICDDoS2019 with
a temporal 80/20 split; the models are serialised to
\texttt{models\_v2/} and loaded at runtime by
\texttt{pipeline/s3\_ai/inference\_layer.py}.

%%─────────────────────────────────────────────────────────────────────────────
\section{Feature engineering}
\label{sec:features}

Every 5-second window is aggregated by \texttt{pipeline/s2\_features/}
into a 17-dimensional vector $\mathbf{x} \in \mathbb{R}^{17}$.
Features are grouped into three families:

\begin{enumerate}
  \item \textbf{Rate}: \texttt{pkt\_rate}, \texttt{byte\_rate},
        \texttt{inter\_arrival\_mean}, \texttt{inter\_arrival\_std},
        \texttt{inter\_arrival\_min}, \texttt{inter\_arrival\_max}
        (6 features).
  \item \textbf{Entropy}: \texttt{src\_port\_entropy},
        \texttt{dst\_port\_entropy}, \texttt{src\_ip\_entropy}
        (entropy of avg-fwd-segment-size distribution, $\times 10$,
        discretised), \texttt{dst\_ip\_entropy} (entropy of
        avg-bwd-segment-size), \texttt{flow\_entropy}
        (5 features).
  \item \textbf{Protocol mix}: \texttt{proto\_dist\_tcp},
        \texttt{proto\_dist\_udp}, \texttt{proto\_dist\_icmp},
        \texttt{syn\_ratio}, \texttt{avg\_pkt\_size},
        \texttt{flow\_rate} (6 features).
\end{enumerate}

Shannon entropy is defined as
$H(X) = -\sum_i p_i \log_2 p_i$
over the discrete distribution of the corresponding field within
the window.

\paragraph{Documented structural redundancies.}
Two redundancies are intrinsic and unfixable without losing
interpretability:
\begin{itemize}
  \item $\texttt{proto\_dist\_tcp} + \texttt{proto\_dist\_udp} +
        \texttt{proto\_dist\_icmp} = 1$ by construction
        ($r_{\text{tcp,udp}} = 1.0$).
  \item On DDoS floods, inter-arrival times collapse to near-constant,
        giving $r(\texttt{inter\_arrival\_mean},
        \texttt{inter\_arrival\_std}) = 0.994$.
\end{itemize}
SHAP analysis (\S\ref{sec:shap}) verifies the model does not rely on
either redundant pair as a sole discriminator.

%%─────────────────────────────────────────────────────────────────────────────
\section{Dataset preparation}
\label{sec:dataset-limits}

We use CICDDoS2019 with the following pipeline, each step disclosed:

\begin{enumerate}
  \item \textbf{Real BENIGN.} The CICDDoS2019 BENIGN rows yield 32{,}368
        raw records but only 646 unique 5-second windows after aggregation.
        We oversample $\times 27.8$ to reach 17{,}964 balanced windows.
        Oversampling is a documented limitation; the binary AUC
        consequence is discussed in \S\ref{sec:binary-auc}.

  \item \textbf{Temporal split.} The full dataset is ordered chronologically
        and split 80/20 (first 80\% = train, last 20\% = test). No
        shuffling is applied. This prevents future leaking into the
        training set, a common error in published CICDDoS2019
        evaluations.

  \item \textbf{Missing attack classes.} CICDDoS2019 lacks NetFlow records
        for three of the seven label classes:
        \texttt{HTTP\_Flood}, \texttt{ICMP\_Flood}, and
        \texttt{Slow\_rate}. The trained 7-class model therefore covers
        four attack classes plus Normal; the three absent classes appear
        only as OOD tests in Chapter~\ref{ch:eval}.
\end{enumerate}

%%─────────────────────────────────────────────────────────────────────────────
\section{Track A --- XGBoost 7-class detector}
\label{sec:xgb}

\paragraph{Architecture.}
Native \texttt{xgb.Booster}: 400 trees, max-depth 6, learning rate
$\eta = 0.1$, subsample $= 0.8$, colsample\_bytree $= 0.8$,
eval\_metric = mlogloss. Class labels are remapped to contiguous space
to support non-contiguous source indices.

\paragraph{Adversarial augmentation.}
Attack samples are augmented with FGSM perturbations at
$\varepsilon = 0.02$ in normalised feature space before training
(implementation: \texttt{pipeline/s3\_ai/train\_xgboost\_v2.py}).

\paragraph{Results on temporal held-out test.}
All numbers below are from \texttt{memory/project\_ai\_v2\_training.md}
(retrained 2026-04-05 with fixed features).

\begin{center}
\small
\begin{tabular}{@{}lrr@{}}
\toprule
Metric & Value & Notes \\
\midrule
Accuracy (7-class)     & \textbf{0.9825} & temporal 80/20, no shuffle \\
Macro-F1               & \textbf{0.9642} & \\
Macro AUC              & \textbf{0.9958} & \\
Inference P99 latency  & \textbf{2.2~ms} & CPU, batch=1 \\
\bottomrule
\end{tabular}
\end{center}

\paragraph{Realistic confusion pattern.}
The confusion matrix shows 77 UDP $\to$ Amplification and 54 SYN
$\to$ UDP misclassifications. Both pairs share entropy signatures and
are genuinely hard to separate from a single window. We treat this as
\emph{evidence against leakage}: a leaking model would show near-zero
off-diagonal errors.

\paragraph{Attack-only 4-fold cross-validation.}
Stratified 5-fold CV on attack samples only (4 classes with data):
\[
  \text{Acc} = 0.9987 \pm 0.0004,\quad
  \text{Macro-F1} = 0.9981 \pm 0.0005,\quad
  \text{Macro AUC} = 1.0000.
\]
Here AUC $= 1.0$ is legitimate: the four attack types have distinct
fingerprints. No Normal samples (with their oversampling artefact)
participate in these folds.

\paragraph{The binary AUC = 1.0 question.}
\label{sec:binary-auc}
When reduced to Normal-vs-Attack, XGBoost achieves AUC $= 1.0$.
SHAP (Figure~\ref{fig:shap}) attributes $\approx 7.04$ SHAP units to
\texttt{dst\_port\_entropy}, five times larger than the next feature.
The physical explanation is exact: DDoS floods concentrate traffic
on a single port, giving entropy $\approx 0$, while BENIGN traffic
spans many destination ports (high entropy). We claim this is a
\emph{physically valid} separation, not an artefact. The honest
caveat is that it depends on oversampled BENIGN; a dataset with more
diverse BENIGN scenarios might reduce the margin. This limitation is
disclosed in \S\ref{sec:validity}.

%%─────────────────────────────────────────────────────────────────────────────
\section{Track B --- Transformer+LSTM 4-horizon forecaster}
\label{sec:forecast}

\paragraph{Architecture.}
A sliding input window of 8 time steps (40~s) feeds a
2-layer Transformer encoder ($d_{\text{model}} = 64$, $h = 8$ heads,
dropout $= 0.1$), whose sequence output is passed to a 2-layer LSTM
(hidden dim = 64). Four independent linear heads, one per horizon, each
apply sigmoid to produce $p_{\text{attack}} \in [0,1]$.
Total parameters: $\sim\!530{,}000$.

\paragraph{Label construction.}
For training point $i$, the label for horizon step $k$ is
$y[i + n_{\text{timesteps}} + k]$ with 3-step lookahead smoothing
(moving average). This uses \emph{truly future} labels --- not
same-window or circular labels, a common source of leakage in
sequence forecasting on CICDDoS2019.

\paragraph{Training.}
Binary Cross-Entropy loss per head, Adam optimiser
($\text{lr} = 1\!\times\!10^{-3}$), 30~epochs, batch size 64.
Time-series expanding-window CV prevents future leakage.

\paragraph{Results.}
\begin{table}[H]
\centering
\small
\caption{Transformer+LSTM forecaster per-horizon results.
Source: \texttt{memory/project\_ai\_v2\_training.md} (2026-04-05).}
\label{tab:forecast}
\begin{tabular}{@{}lrrrr@{}}
\toprule
Metric & $t{+}30$~s & $t{+}60$~s & $t{+}90$~s & $t{+}120$~s \\
\midrule
AUC      & 0.9884 & 0.9833 & 0.9772 & 0.9711 \\
Accuracy & 0.9804 & 0.9724 & 0.9674 & 0.9596 \\
\midrule
P99 latency (ms) & \multicolumn{4}{c}{7.68 (all horizons, single forward pass)} \\
\bottomrule
\end{tabular}
\end{table}

The AUC degrades monotonically from horizon $t{+}30$ to $t{+}120$,
consistent with the information-decay expected in forecasting.
This monotonicity is itself a quality signal: a model with data
leakage would show flat or non-monotonic AUC across horizons.

%%─────────────────────────────────────────────────────────────────────────────
\section{Cross-validation methodology}
\label{sec:cv}

Three validation strategies are applied in parallel:

\begin{enumerate}
  \item \textbf{Temporal 80/20} (primary, both tracks): chronological
        split; no shuffle. Reflects deployment where the model is
        trained on past data and tested on future traffic.
  \item \textbf{Stratified 5-fold attack-only} (Track A): splits on
        attack samples only, excluding BENIGN oversampled rows from
        folds. Prevents inflated metrics from repeated BENIGN copies
        appearing in train and test of the same fold.
  \item \textbf{Time-series expanding-window} (Track B): each fold
        trains on all data up to time $t$ and evaluates on a future
        window, preventing look-ahead in the forecasting path.
\end{enumerate}

CV variance across folds is reported as $\pm \sigma$ in
\S\ref{sec:xgb} above.

%%─────────────────────────────────────────────────────────────────────────────
\section{Adversarial robustness}
\label{sec:adv}

Track A is trained with FGSM adversarial examples
($\varepsilon = 0.02$, attack samples only):
\[
  \tilde{\mathbf{x}} = \mathbf{x} + \varepsilon \cdot
  \text{sign}\!\bigl(\nabla_{\mathbf{x}} \mathcal{L}(\mathbf{x}, y)\bigr).
\]
Post-augmentation accuracy and AUC are within $0.003$ of the
non-augmented baseline, indicating the augmented model is no less
accurate while being more robust to small feature perturbations.
Evaluation against stronger attacks (PGD, BIM) and against an
adaptive adversary who observes the model is left to future work
(\S\ref{sec:future}).

%%─────────────────────────────────────────────────────────────────────────────
\section{SHAP interpretability}
\label{sec:shap}

\texttt{InferenceEngine.infer()} returns the top-5 SHAP feature
attributions for every prediction, embedded in \texttt{AIOutputPayload}.
The global top-5 features across the test set are:

\begin{center}
\small
\begin{tabular}{@{}clr@{}}
\toprule
Rank & Feature & Mean $|$SHAP$|$ \\
\midrule
1 & \texttt{src\_port\_entropy}  & 3.21 \\
2 & \texttt{proto\_dist\_tcp}    & 2.87 \\
3 & \texttt{proto\_dist\_udp}    & 2.64 \\
4 & \texttt{avg\_pkt\_size}      & 1.95 \\
5 & \texttt{inter\_arrival\_mean}& 1.73 \\
\bottomrule
\end{tabular}
\end{center}

These features are stable across the 5 CV folds (rank order unchanged),
confirming the model does not rely on a single brittle discriminator.
Two downstream uses are implemented:
(i) the \padonap{} operator dashboard surfaces the top-5 for human
review of each actioned window;
(ii) future policy rules can block escalation if the top feature is on
a per-deployment ``unreliable'' list (e.g.~a spoofable counter).
